\documentclass[11pt, letterpaper]{article}
\usepackage{amsmath, amssymb}
\usepackage{hyperref}
\usepackage{booktabs}
\usepackage{enumitem}
\usepackage{xcolor}
\usepackage{geometry}
\usepackage{fancyhdr}
\usepackage{titlesec}
\usepackage{tcolorbox}
\usepackage{graphicx}
\usepackage{parskip}
\usepackage{caption}
\geometry{margin=1in}

\definecolor{courseblue}{HTML}{002452}
\definecolor{coursegold}{HTML}{FABD0F}
\definecolor{coursered}{HTML}{B90E31}
\definecolor{lightblue}{HTML}{E8EEF7}
\definecolor{lightgold}{HTML}{FEF8E3}

\hypersetup{colorlinks=true,linkcolor=courseblue,urlcolor=coursered}
\tcbuselibrary{skins}
\newtcolorbox{guidancebox}{enhanced,colback=lightblue,colframe=courseblue,
  sharp corners,leftrule=4pt,left=6pt,right=6pt,top=4pt,bottom=4pt}
\newtcolorbox{notebox}{enhanced,colback=lightgold,colframe=coursegold,
  sharp corners,leftrule=4pt,left=6pt,right=6pt,top=4pt,bottom=4pt}

\pagestyle{fancy}\fancyhf{}
\fancyhead[L]{\small\textcolor{courseblue}{\textbf{CISC 867 --- Deep Learning | Fall 2026}}}
\fancyhead[R]{\small\textcolor{courseblue}{Midterm Report}}
\fancyfoot[C]{\small\thepage}
\renewcommand{\headrulewidth}{0.4pt}
\titleformat{\section}{\large\bfseries\color{courseblue}}{}{0em}{}[\titlerule]
\titleformat{\subsection}{\normalsize\bfseries\color{courseblue}}{}{0em}{}

\begin{document}

{\centering
  {\Large\textbf{CISC 867: Deep Learning --- Midterm Report}}\\[0.5em]
  {\large\textit{[PROJECT ID AND TITLE]}}\\[0.4em]
  \textbf{Group Members:}\\
  {[Student 1 Name]} (Project Lead) \quad
  {[Student 2 Name]} (Research Lead) \quad
  {[Student 3 Name]} (Engineering Lead)\\[0.3em]
  \textbf{GitHub:} \url{[link]} \quad \textbf{Date:} \today\\[0.4em]
  \rule{\linewidth}{0.4pt}
}

\section{Abstract}

\begin{guidancebox}
\textbf{Guidance.} 150--200 word summary of: (1) the problem, (2) what you have
implemented so far, (3) your key preliminary result, and (4) your planned next steps.
\end{guidancebox}

Diffusion models can generate high-quality images, but conventional Denoising Diffusion Probabilistic Models (DDPMs) require many sequential denoising steps, making image generation computationally expensive. This project investigates the trade-off between image quality and sampling efficiency on MNIST and CIFAR-10, with the longer-term goal of combining quality, generation speed, and computational cost in a single evaluation metric. As a baseline, a DDPM was implemented from scratch using a compact U-Net noise predictor, sinusoidal timestep embeddings, a 1,000-step linear noise schedule, and a mean squared error noise-prediction objective. The saved MNIST run completed 20 training epochs and reports a preliminary Fr\'echet Inception Distance (FID) of 5.72. This is an initial result, not a controlled comparison; preprocessing and evaluation settings will be checked before drawing conclusions. Next, the model will be ported to CIFAR-10 and compare linear and cosine schedules, implement DDIM and measure FID at 10, 50, and 1,000 sampling steps, and evaluate classifier-free guidance and a single-step Consistency Model sampler. Sampling-time and computational-cost measurements will support the proposed quality-efficiency metric and comparisons.


\section{Introduction and Problem Statement}

\begin{guidancebox}
\textbf{Guidance.} Brief recap of the problem and why it matters (1/2 page max).
Focus on what has changed since the proposal---have you refined the scope or identified
a clearer novel angle? Include the \textbf{research question} you are now testing.
\end{guidancebox}

[CONCISE PROBLEM DESCRIPTION AND RESEARCH QUESTION. NOTE ANY SCOPE CHANGES.]

Diffusion models generate images by learning to reverse gradual noising of training examples. A conventional Denoising Diffusion Probabilistic Model (DDPM) samples through many sequential denoising steps, making generation computationally costly. Denoising Diffusion Implicit Models (DDIMs) and Consistency Models aim to reduce the number of model evaluations, but image quality and generation cost must be considered together. FID measures similarity between generated and real image distributions, but does not capture inference time or computational cost.

The proposal's broad aim was to compare diffusion methods and develop a metric combining quality and efficiency. The project guide specifies staged tests on MNIST and CIFAR-10: establish a DDPM baseline, compare linear and cosine noise schedules, measure DDIM FID at 10, 50, and 1,000 steps, and assess classifier-free guidance and a single-step Consistency Model sampler. So far, a saved MNIST DDPM run reports an FID of 5.72 after 20 epochs. This is an initial result, not a controlled comparison; preprocessing and evaluation still require verification. These experiments will investigate whether faster sampling can reduce computational cost while preserving image quality, and whether a single metric can clearly summarize that trade-off.

\section{Related Work (Updated)}

\begin{guidancebox}
\textbf{Guidance.} An extended version of your proposal survey. Add papers you
discovered since the proposal. Include a comparative table or taxonomy where possible.
Make explicit connections to your own approach. Add at least 2 new papers beyond
the proposal. Target: $\approx$400--600 words.
\end{guidancebox}


The proposal introduced DDPMs and faster sampling methods as ways to study the balance between generated-image quality and computational cost. This updated discussion examines two additional methods relevant to that goal: DPM-Solver, a faster sampler for an existing diffusion model, and Classifier-Free Diffusion Guidance, which steers generation toward a requested class. Neither method has been implemented in this project. They offer complementary directions for future work: reducing the number of sampling evaluations and enabling class-conditional generation.

\paragraph{\textit{DPM-Solver: A Fast ODE Solver for Diffusion Probabilistic Model Sampling in Around 10 Steps} (Lu et al., NeurIPS 2022).}
Lu et al. treat diffusion sampling as solving a probability-flow ordinary differential equation and introduce a solver designed to approximate its solution with relatively few evaluations of the diffusion model. Unlike methods that require training a new generative model, DPM-Solver is a sampling procedure intended to work with an existing diffusion model. This makes it relevant to our efficiency investigation: it could potentially be applied to the trained MNIST denoiser and compared with the standard DDPM reverse process using FID and sampling cost at different numbers of model evaluations. We have not implemented DPM-Solver; its compatibility with our discrete-time model and available project time would need to be checked before adding it to the experiments.

\paragraph{\textit{Classifier-Free Diffusion Guidance} (Ho and Salimans, 2022).}
Ho and Salimans introduce a way to guide conditional diffusion generation without training a separate image classifier. During training, the conditioning information is sometimes withheld, teaching the diffusion network both conditional and unconditional noise predictions. At sampling time, those predictions are combined using a guidance scale to steer generation toward the requested condition. Stronger guidance can improve alignment with the condition, but may affect sample diversity and quality. This method is relevant to the project guide's objective of testing class-conditional generation on MNIST. Our current U-Net ignores digit labels, so applying this method would require adding label conditioning and training with label dropout. We have not implemented this change; the planned experiment is to compare guided and unguided class-conditional samples and assess their quality.

\begin{center}
\begin{tabular}{llp{4.5cm}p{5cm}}
\toprule
\textbf{Method} & \textbf{Year} & \textbf{Key idea} & \textbf{Relation to our work} \\
\midrule
DPM-Solver & 2022 & Solve the probability-flow ODE with a method designed to reduce diffusion-model evaluations. & Not implemented. We could test it as an alternative sampler for the trained denoiser, comparing quality and sampling cost if compatible. \\
\midrule
Classifier-Free Guidance & 2022 & Train with and without conditioning, then combine both predictions at sampling time to guide generation without a separate classifier. & Not implemented. We plan to condition the MNIST model on digit labels and compare guided with unguided generation. \\
\bottomrule
\end{tabular}
\end{center}

Together, these methods motivate complementary experiments rather than interchangeable model comparisons. DPM-Solver addresses sampling efficiency by reducing model evaluations, while classifier-free guidance adds control over the generated digit class. Any future comparisons would need consistent datasets and evaluation procedures; published FID scores from different experiments should not be treated as directly comparable.



\section{Methodology}

\begin{guidancebox}
\textbf{Guidance.} Describe what you have implemented in detail. Include:
(1) architecture / model description with any modifications,
(2) training procedure, hyperparameters, and optimiser,
(3) data preprocessing pipeline,
(4) evaluation protocol (metric definitions and dataset splits).
Use equations where appropriate.
\end{guidancebox}

\subsection{Model Architecture}

[DESCRIBE YOUR MODEL. If you are using a pre-trained backbone, state clearly
which weights you use and what you add/modify on top.]

We implemented and trained a denoising diffusion model to generate handwritten digits from the MNIST dataset. It uses a U-Net neural network to predict and remove noise from images, starting from random noise; the model was trained from randomly initialized weights rather than pretrained weights. The network takes a one-channel, $28\times28$ noisy image and its diffusion timestep as inputs, and predicts the Gaussian noise added to the image. The timestep is encoded with a 128-dimensional sinusoidal embedding and a two-layer multilayer perceptron with a SiLU activation. The image path begins with a convolution that maps the input to 64 channels. Two encoder residual blocks process features at successively lower spatial resolutions, with stride-2 convolutions increasing the channel widths from 64 to 128 and then 256. Two residual blocks form the bottleneck. Transposed convolutions then upsample the features; at each decoder stage, the features are concatenated with the matching encoder features through a skip connection. The decoder returns a one-channel noise prediction. Residual blocks use Group Normalization with eight groups, SiLU activations, $3\times3$ convolutions, timestep conditioning, and identity or $1\times1$ convolutional residual paths as required by the channel dimensions. The model has 4,788,417 trainable parameters.

\subsection{Training Setup}

For each image in a training batch, let $x_0$ denote the original, clean image. The notebook independently chooses a timestep $t$ uniformly from the $T=1{,}000$ forward-diffusion steps and draws Gaussian noise $\epsilon\sim\mathcal{N}(0,I)$ with the same shape as the image. Here, $I$ is the identity covariance matrix, so the noise values are standard normal. At step $t$, the noise schedule specifies $\beta_t$, the variance of the noise added at that step; $\alpha_t=1-\beta_t$ is the fraction of signal retained at that step, and $\bar{\alpha}_t=\prod_{s=1}^{t}\alpha_s$ is the cumulative signal retained through step $t$. The resulting noisy image is
\[
x_t=\sqrt{\bar{\alpha}_t}\,x_0+\sqrt{1-\bar{\alpha}_t}\,\epsilon,
\]
where $x_t$ is the image after adding noise up to timestep $t$. The U-Net, denoted by $\epsilon_\theta$, receives $x_t$ and $t$ and predicts the noise $\epsilon$ that was added. Training minimizes the mean squared error between that prediction and the known noise:
\[
\mathcal{L}=\mathbb{E}_{x_0,t,\epsilon}\!\left[\operatorname{MSE}\!\left(\epsilon,\epsilon_\theta(x_t,t)\right)\right].
\]
The expectation indicates that this error is averaged over training images, randomly selected timesteps, and sampled noise. The saved run uses AdamW and completed 20 epochs on a CUDA-enabled Colab T4 GPU. To generate images, the sampler begins with Gaussian noise and repeatedly applies the learned reverse update from timestep $T$ down to 1, adding noise at each step except the final one.
\newline

\begin{center}
\begin{tabular}{ll}
\toprule
\textbf{Hyperparameter} & \textbf{Value} \\
\midrule
Optimiser & AdamW \\
Learning rate & $2\times10^{-4}$ \\
Weight decay & $1\times10^{-4}$ \\
Batch size & 128 \\
Epochs & 20 \\
Diffusion steps & 1,000 \\
Noise schedule & Linear, $\beta_{\min}=10^{-4}$ to $\beta_{\max}=0.02$ \\
Loss function & Mean squared error on predicted Gaussian noise \\
Data augmentation & None \\
Random seeds & Not fixed or reported; one saved run \\
Trainable parameters & 4,788,417 \\
\bottomrule
\end{tabular}
\end{center}

\subsection{Dataset and Preprocessing}

[DESCRIBE THE DATASET SPLIT (train/val/test sizes), ANY PREPROCESSING STEPS
(resizing, normalisation, augmentation), AND HOW YOU HANDLE CLASS IMBALANCE
IF APPLICABLE.]

To train the handwritten-digit generator, we used the 60,000-image MNIST training split, loaded with \texttt{torchvision.datasets.MNIST(train=True)}. The images were shuffled into batches of 128, and all training examples were used in each of the 20 epochs. We did not create a validation split or load the official 10,000-image MNIST test split; the model was therefore not selected or evaluated on held-out test data. For the preliminary FID calculation, generated samples were compared with images exported from the training split, rather than a separate test set. Digit labels were ignored, so the model generates digits unconditionally instead of generating a requested class. We converted images to tensors with \texttt{ToTensor()}, scaling pixel values to $[0,1]$. We applied no resizing, additional normalization, augmentation, or class-balancing procedure. CIFAR-10 will be used in planned experiments but was not included in this training run.

\subsection{Evaluation Protocol}

The notebook evaluates sample quality with FID, computed by \texttt{pytorch-fid} using 2,048-dimensional Inception features and a batch size of 50. FID compares the feature distributions of real and generated images; lower scores indicate distributions that are closer under this measure. The saved run generates 10,000 MNIST samples using the full 1,000-step reverse process and compares them with images exported from the MNIST training split. The real-image export loop saves 10,001 images, converts them to RGB by repeating the grayscale channel, and the generated images are likewise converted to RGB before scoring. 

The preprocessing in this FID calculation needs to be corrected or verified before the score is treated as a reliable baseline: training images enter the model in $[0,1]$, but the real-image export applies $(x+1)/2$, whereas generated outputs are clamped to $[-1,1]$ and mapped to $[0,1]$. Thus, real and generated images are not exported using matching intensity transformations. The evaluation also uses training images as its real reference set rather than a held-out test set. These limitations will be addressed when establishing a reproducible comparison protocol.


\section{Preliminary Results}

\begin{guidancebox}
\textbf{Guidance.} Present quantitative results for your baseline (and, if
available, early results for your proposed method). Report mean $\pm$ std over
at least 3 seeds. Compare to at least one published baseline. Include at least
one qualitative example (figure or table). Discuss what the numbers tell you.
\end{guidancebox}

\subsection{Baseline Reproduction}

\begin{notebox}
\textbf{Target:} Match the published number within 1--2\% before proceeding
to your proposed modification.
\end{notebox}

\begin{center}
\begin{tabular}{lcc}
\toprule
\textbf{Model} & \textbf{Metric 1 (Mean $\pm$ Std)} & \textbf{Metric 2 (Mean $\pm$ Std)} \\
\midrule
Published baseline & [Value from paper] & [Value from paper] \\
Our reproduction & [Your value] $\pm$ [Std] & [Your value] $\pm$ [Std] \\
Our proposed (early) & [Your value] $\pm$ [Std] & [Your value] $\pm$ [Std] \\
\bottomrule
\end{tabular}
\captionof{table}{Preliminary results on [DATASET], [SPLIT]. All results are mean $\pm$ std over 3 seeds.}
\end{center}

\subsection{Qualitative Analysis}

[INSERT A REPRESENTATIVE FIGURE: e.g., a before/after image pair, a learning
curve, a confusion matrix, or a prediction example. Caption should fully
describe what is being shown.]

\begin{figure}[h]
  \centering
  % \includegraphics[width=0.8\linewidth]{figures/qualitative.png}
  \caption{[FIGURE CAPTION: describe the figure and what it shows.]}
  \label{fig:qual}
\end{figure}

\subsection{Discussion}

[INTERPRET YOUR RESULTS. What works? What does not? What surprised you?
Where are the biggest failure cases? What do you conclude so far?]

\section{Challenges and Adaptations}

\begin{guidancebox}
\textbf{Guidance.} Describe any obstacles you encountered and how you
addressed or plan to address them. Did you deviate from your proposal timeline?
If so, explain why and provide an updated plan.
\end{guidancebox}

[DESCRIBE: (1) Technical obstacles (bugs, slow training, dataset issues),
(2) Conceptual challenges (method harder than expected, results not matching),
(3) How you resolved or plan to resolve each issue,
(4) Updated timeline for Weeks 7--10.]

\subsection{Updated Week 7--10 Plan}

\begin{tabular}{lp{8cm}l}
\toprule
\textbf{Week} & \textbf{Tasks} & \textbf{Responsible} \\
\midrule
7 & [Updated tasks] & [Name] \\
8 & [Updated tasks] & [Name] \\
9 & [Updated tasks] & [Name] \\
10 & [Final report + presentation] \textbf{(Final due)} & All \\
\bottomrule
\end{tabular}

\section{Preliminary Novel Contribution (Optional Bonus)}

\begin{guidancebox}
\textbf{Guidance.} If you have a preliminary finding that goes beyond
reproducing the baseline---an unexpected result, an improved ablation,
an insight about the method---describe it here for the +2 bonus.
Provide quantitative evidence.
\end{guidancebox}

[DESCRIBE YOUR PRELIMINARY NOVEL FINDING, WITH EVIDENCE. OR LEAVE EMPTY.]

\section{Contribution Declaration}

\begin{tabular}{lcc}
\toprule
\textbf{Member} & \textbf{Contributions (this period)} & \textbf{\%} \\
\midrule
{[Student 1]} & [e.g., Lit review, updated proposal section] & \% \\
{[Student 2]} & [e.g., Model implementation, training runs] & \% \\
{[Student 3] }& [e.g., Data pipeline, evaluation code] & \% \\
\bottomrule
\end{tabular}

\section{References}

\begin{enumerate}[label={[\arabic*]},leftmargin=*]
  \item Cheng Lu, Yuhao Zhou, Fan Bao, Jianfei Chen, Chongxuan Li, Jun Zhu, 2022. ``DPM-Solver: A Fast ODE Solver for Diffusion Probabilistic Model Sampling in Around 10 Steps.'' \textit{Advances in Neural Information Processing Systems}. arXiv:2206.00927
  \item Jonathan Ho and Tim Salimans, 2022. ``Classifier-Free Diffusion Guidance.'' arXiv:2207.12598
\end{enumerate}
\end{enumerate}

\end{document}
